\documentclass[12pt, a4paper]{article}
\usepackage[utf8]{inputenc}
\usepackage[T1]{fontenc}
\usepackage[ngerman]{babel}
\usepackage{amsmath, amssymb, amsthm}
\usepackage{geometry}
\usepackage{graphicx}
\usepackage{xcolor}
\usepackage{tcolorbox}
\usepackage{enumitem}
\usepackage{listings}
\usepackage{fancyhdr}
\usepackage{titlesec}
\usepackage{mdframed}
\usepackage{tikz}
\usepackage{pgfplots}
\pgfplotsset{compat=1.18}

\geometry{a4paper, margin=2.5cm}

\definecolor{theoryblue}{RGB}{30, 90, 160}
\definecolor{exerciseorange}{RGB}{200, 80, 0}
\definecolor{solutiongreen}{RGB}{0, 120, 60}
\definecolor{hintgray}{RGB}{100, 100, 100}
\definecolor{codebg}{RGB}{245, 245, 245}

% Theory box
\tcbuselibrary{skins, breakable}
\newtcolorbox{theorybox}[1]{
  colback=theoryblue!5,
  colframe=theoryblue,
  fonttitle=\bfseries\color{white},
  title=#1,
  breakable
}

% Exercise box
\newtcolorbox{exercisebox}[1]{
  colback=exerciseorange!5,
  colframe=exerciseorange,
  fonttitle=\bfseries\color{white},
  title=Aufgabe #1,
  breakable
}

% Solution space box
\newtcolorbox{solutionspace}[1][6cm]{
  colback=white,
  colframe=black!20,
  fonttitle=\itshape\color{hintgray},
  title=Rechenraum,
  height=#1,
  breakable
}

% Hint box
\newtcolorbox{hintbox}{
  colback=solutiongreen!5,
  colframe=solutiongreen,
  fonttitle=\bfseries\color{white},
  title=Hinweis,
  breakable
}

% Julia code style
\lstdefinelanguage{Julia}{
  keywords={function, end, return, if, else, elseif, for, while, in, begin, let, do, try, catch, finally, using, import, export, mutable, struct, abstract, type},
  keywordstyle=\color{theoryblue}\bfseries,
  comment=[l]{\#},
  commentstyle=\color{hintgray}\itshape,
  stringstyle=\color{solutiongreen},
  morestring=[b]",
  sensitive=true
}
\lstset{
  language=Julia,
  backgroundcolor=\color{codebg},
  basicstyle=\ttfamily\small,
  breaklines=true,
  frame=single,
  rulecolor=\color{black!30},
  numbers=left,
  numberstyle=\tiny\color{hintgray}
}

\pagestyle{fancy}
\fancyhf{}
\rhead{\textcolor{hintgray}{\small Rocket Science Workbook}}
\lhead{\textcolor{hintgray}{\small \leftmark}}
\cfoot{\thepage}

\titleformat{\section}{\large\bfseries\color{theoryblue}}{}{0em}{}[\titlerule]
\titleformat{\subsection}{\normalsize\bfseries\color{exerciseorange}}{}{0em}{}


\begin{document}

\begin{center}
  {\LARGE \textbf{Chapter 02: Runge-Kutta 4th Order (RK4)}}\\[0.5em]
  {\large \textcolor{hintgray}{From Simpson's Rule to the Most Important ODE Solver in Science}}\\[0.3em]
  {\small \textcolor{hintgray}{Rocket Science Workbook — simulation-2D-jl}}
\end{center}
\vspace{1em}\hrule\vspace{1.5em}

% ──────────────────────────────────────────────────────────────
\section{The Integral Form of an ODE}

\begin{theorybox}{From Derivative to Integral}
The ODE $\dot{y} = f(t,y)$ can be integrated exactly over $[t_n, t_{n+1}]$:
\[
  y(t_{n+1}) = y(t_n) + \int_{t_n}^{t_n+h} f(\tau,\, y(\tau))\,d\tau
\]
The problem: we don't know $y(\tau)$ inside the integral — that's what we're trying to find.\\[0.4em]
Every numerical method is a \textbf{quadrature rule} for this integral.
\end{theorybox}

\subsection*{Three Classical Quadrature Rules}

\begin{center}
\begin{tabular}{lll}
\textbf{Method} & \textbf{Approximation of integral} & \textbf{Error} \\
\hline
Rectangle (left) & $h\cdot f(t_n, y_n)$ & $\mathcal{O}(h^2)$ $\to$ Euler \\
Midpoint rule & $h\cdot f(t_n+h/2,\; y_n + \frac{h}{2}f_n)$ & $\mathcal{O}(h^3)$ $\to$ RK2 \\
Simpson's rule & $\frac{h}{6}(f_a + 4f_m + f_b)$ & $\mathcal{O}(h^5)$ $\to$ RK4 \\
\end{tabular}
\end{center}

RK4 is essentially \textbf{Simpson's rule} applied adaptively to the ODE integral.

% ──────────────────────────────────────────────────────────────
\section{Deriving RK4 from Taylor Series}

\begin{theorybox}{Why Euler fails at order 1}
From Chapter 01: $y(t+h) = y + h f + \frac{h^2}{2}(f_t + f_y f) + \mathcal{O}(h^3)$

Euler gives only $y + hf$ — the $\mathcal{O}(h^2)$ term is the error.
To cancel it without computing $\partial f/\partial t$ and $\partial f/\partial y$ explicitly,
we can \textbf{evaluate $f$ at multiple cleverly chosen points} and combine them.
\end{theorybox}

\begin{theorybox}{RK4 — The Four Stages}
\[
  k_1 = f(t_n,\; y_n)
\]
\[
  k_2 = f\!\left(t_n + \tfrac{h}{2},\; y_n + \tfrac{h}{2}k_1\right)
\]
\[
  k_3 = f\!\left(t_n + \tfrac{h}{2},\; y_n + \tfrac{h}{2}k_2\right)
\]
\[
  k_4 = f\!\left(t_n + h,\; y_n + h\,k_3\right)
\]
\[
  \boxed{y_{n+1} = y_n + \frac{h}{6}(k_1 + 2k_2 + 2k_3 + k_4)}
\]
\end{theorybox}

\subsection*{Why the Weights $1:2:2:1$?}

Expand each $k_i$ as a Taylor series around $(t_n, y_n)$ and collect terms in powers of $h$.
The weights $1/6,\; 2/6,\; 2/6,\; 1/6$ are chosen exactly so that all terms up to $h^4$ cancel with the Taylor series of the exact solution. The first uncancelled term is $\mathcal{O}(h^5)$.

\begin{theorybox}{Error Orders}
\begin{itemize}
  \item \textbf{Local truncation error:} $\mathcal{O}(h^5)$ per step
  \item \textbf{Global error:} $\mathcal{O}(h^4)$ — halving $h$ reduces error by factor $2^4 = 16$
  \item \textbf{Function evaluations per step:} 4 (compared to 1 for Euler)
\end{itemize}
\end{theorybox}

% ──────────────────────────────────────────────────────────────
\section{The Butcher Tableau}

All Runge-Kutta methods can be summarized in a compact \textbf{Butcher tableau}:
\[
\begin{array}{c|cccc}
0   &      &     &     &   \\
1/2 & 1/2  &     &     &   \\
1/2 &  0   & 1/2 &     &   \\
1   &  0   &  0  &  1  &   \\
\hline
    & 1/6  & 2/6 & 2/6 & 1/6
\end{array}
\]
The column of the left ($c_i$) gives the time offsets; $a_{ij}$ gives how to combine previous stages; $b_i$ are the final weights.

% ──────────────────────────────────────────────────────────────
\section{Taylor Series Verification of $k_2$}

\begin{theorybox}{Expanding $k_2$ around $(t_n, y_n)$}
\[
  k_2 = f\!\left(t_n + \tfrac{h}{2},\; y_n + \tfrac{h}{2}k_1\right)
\]
Using a 2D Taylor expansion of $f(t+\Delta t, y+\Delta y)$:
\[
  k_2 = f + \frac{h}{2}\frac{\partial f}{\partial t} + \frac{h}{2}f\frac{\partial f}{\partial y} + \mathcal{O}(h^2)
\]
where all partial derivatives are evaluated at $(t_n, y_n)$.
Substituting into the RK4 formula and collecting, one finds the first uncancelled term is $\propto h^5$ — confirming 4th order accuracy.
\end{theorybox}

% ──────────────────────────────────────────────────────────────
\section{Exercises}

\begin{exercisebox}{1 — Simpson's Rule Connection}
Simpson's rule for $\int_a^b g(x)\,dx$ uses three points:
\[
  \int_a^b g(x)\,dx \approx \frac{b-a}{6}\left(g(a) + 4g\!\left(\frac{a+b}{2}\right) + g(b)\right)
\]

\textbf{a)} Write out the RK4 update formula in terms of $k_1, k_2, k_3, k_4$ and identify
the Simpson weights $(1, 4, 1)$ hidden in the $(1, 2, 2, 1)$ pattern.
\textit{Hint: think of $k_2$ and $k_3$ as two estimates of the midpoint.}

\textbf{b)} Show that Simpson's rule is \textit{exact} for polynomials up to degree 3.
\textit{Hint: check for $g(x) = x^3$ on $[0,1]$:}
$\int_0^1 x^3\,dx = 1/4$, \; Simpson gives $\frac{1}{6}(0 + 4\cdot(1/2)^3 + 1) = ?$
\end{exercisebox}

\begin{solutionspace}[8cm]
\end{solutionspace}

\begin{exercisebox}{2 — RK4 by Hand: One Full Step}
$\dfrac{dy}{dt} = -y$, \quad $y(0) = 1$, \quad $h = 0.5$.

Compute all four stages, showing every arithmetic step:

\begin{align*}
k_1 &= f(0.0,\; 1.0000) = -1.0000\\[0.4em]
k_2 &= f\!\left(0.25,\; 1.0 + 0.25\cdot k_1\right) = f(0.25,\; \underline{\hspace{2cm}}) = \underline{\hspace{2cm}}\\[0.4em]
k_3 &= f\!\left(0.25,\; 1.0 + 0.25\cdot k_2\right) = f(0.25,\; \underline{\hspace{2cm}}) = \underline{\hspace{2cm}}\\[0.4em]
k_4 &= f\!\left(0.50,\; 1.0 + 0.50\cdot k_3\right) = f(0.50,\; \underline{\hspace{2cm}}) = \underline{\hspace{2cm}}\\[0.8em]
y_1 &= 1.0 + \frac{0.5}{6}(k_1 + 2k_2 + 2k_3 + k_4) = \underline{\hspace{5cm}}
\end{align*}

Exact: $e^{-0.5} = 0.60653$. \quad
Error RK4: \underline{\hspace{2cm}} \quad
Error Euler ($y_1 = 0.5$): \underline{\hspace{2cm}}

\textbf{How many times more accurate is RK4 than Euler here?}
\end{exercisebox}

\begin{solutionspace}[5cm]
\end{solutionspace}

\begin{exercisebox}{3 — Error Order Experimental Confirmation}
For $\dot{y}=-y$, $y(0)=1$, compute the global error at $t=1$ for multiple step sizes.
Use your hand calculation from Exercise 2 as one data point.

\begin{center}
\begin{tabular}{|c|c|c|c|c|}
\hline
$h$ & $y_N^{\text{RK4}}$ at $t=1$ & Error $|y_N - e^{-1}|$ & Ratio & Expected ratio \\
\hline
0.50 & (from Ex.~2, continue) & & — & — \\
\hline
0.25 & & & & $2^4 = 16$ \\
\hline
0.10 & & & & \\
\hline
\end{tabular}
\end{center}

\textbf{Confirm:} When $h$ is halved, the error drops by a factor of approximately \underline{\hspace{1cm}}.
\end{exercisebox}

\begin{solutionspace}[7cm]
\end{solutionspace}

\begin{exercisebox}{4 — Vector System: One RK4 Step by Hand}
Free fall: $\mathbf{y} = [h, v]^\top$, $f(\mathbf{y}) = [v,\; -9.81]^\top$, $\mathbf{y}_0 = [500, 0]^\top$, $h_\text{step}=1\,\text{s}$.

Compute the full RK4 step. Show vector arithmetic explicitly:

\begin{align*}
\mathbf{k}_1 &= f(0,\; [500, 0]^\top) = [\underline{\hspace{1.5cm}},\; \underline{\hspace{1.5cm}}]^\top \\[0.4em]
\mathbf{y}^* &= [500,0]^\top + 0.5\cdot\mathbf{k}_1 = [\underline{\hspace{2cm}},\; \underline{\hspace{2cm}}]^\top \\[0.4em]
\mathbf{k}_2 &= f(0.5,\; \mathbf{y}^*) = [\underline{\hspace{2cm}},\; \underline{\hspace{2cm}}]^\top \\[0.4em]
\mathbf{k}_3 &= f(0.5,\; [500,0]^\top + 0.5\cdot\mathbf{k}_2) = [\underline{\hspace{2cm}},\; \underline{\hspace{2cm}}]^\top \\[0.4em]
\mathbf{k}_4 &= f(1.0,\; [500,0]^\top + 1.0\cdot\mathbf{k}_3) = [\underline{\hspace{2cm}},\; \underline{\hspace{2cm}}]^\top \\[0.8em]
\mathbf{y}_1 &= [500,0]^\top + \frac{1}{6}(\mathbf{k}_1 + 2\mathbf{k}_2 + 2\mathbf{k}_3 + \mathbf{k}_4)
= [\underline{\hspace{2cm}},\; \underline{\hspace{2cm}}]^\top
\end{align*}

Exact at $t=1$: $h = 500 - 4.905 = 495.095\,\text{m}$, $v = -9.81\,\text{m/s}$.\\
Error RK4: \underline{\hspace{2cm}}m \quad Error Euler: \underline{\hspace{2cm}}m
\end{exercisebox}

\begin{solutionspace}[6cm]
\end{solutionspace}

\begin{exercisebox}{5 — Stability Region (Challenge)}
The stability region of a method is the set of $z = h\lambda$ (complex) for which the method produces bounded solutions.

For Euler: $|1 + z| \leq 1$ \; (disk of radius 1 centered at $-1$).

For RK4, the amplification factor is:
\[
  R(z) = 1 + z + \frac{z^2}{2} + \frac{z^3}{6} + \frac{z^4}{24}
\]
(the first 5 terms of $e^z$!)

\textbf{a)} Show that $R(z) = e^z + \mathcal{O}(z^5)$, confirming 4th order accuracy.

\textbf{b)} For real $z < 0$ (i.e.\ $\lambda < 0$, stable ODE): find the range of $z$ where $|R(z)| \leq 1$.
\textit{Hint: Try $z = -1, -2, -2.5, -2.8, -3$. Compute $R(z)$ for each.}
The stability boundary for RK4 is at $z \approx -2.785$.

\textbf{c)} For the rocket (typical $\lambda \approx -10\,\text{s}^{-1}$ from drag/mass), what is the maximum allowed step size?
\end{exercisebox}

\begin{solutionspace}[10cm]
\end{solutionspace}

% ──────────────────────────────────────────────────────────────
\section{Resources}

\begin{hintbox}
\textbf{Videos:}
\begin{itemize}
  \item \textbf{3Blue1Brown} — ``Simpson's rule, compound quadrature'':\\
        \texttt{youtube.com/watch?v=QXumu31oVdU} (intuition for integration rules)
  \item \textbf{The Coding Train} — ``Runge-Kutta Integration'' (visual, code-focused):\\
        Search: ``Coding Train Runge Kutta''
  \item \textbf{MIT 18.330} — Numerical methods lectures on YouTube (rigorous derivation)
\end{itemize}
\textbf{Reading:}
\begin{itemize}
  \item Burden \& Faires, \textit{Numerical Analysis} Ch.~5 (standard reference)
  \item \texttt{mathworld.wolfram.com/Runge-KuttaMethod.html}
\end{itemize}
\end{hintbox}

\section{Connection to the Julia Template}

\begin{theorybox}{What you implement in \texttt{template.jl}}
\begin{enumerate}
  \item \texttt{rk4\_step(f, t, y, h)} — must work for scalars AND vectors
  \item \texttt{rk4\_solve(f, y0, t0, t\_end, h)} — full integration
  \item Log-scale error plot: Euler $\mathcal{O}(h)$ vs RK4 $\mathcal{O}(h^4)$ — visually striking
  \item Vector system test (Exercise 4)
  \item Stability experiment: solve $\dot{y}=-10y$ with RK4 for $h=0.3$ vs $h=0.28$ — observe blowup
\end{enumerate}
\end{theorybox}

\end{document}
